\documentclass[11pt,letterpaper]{article}
\usepackage[T1]{fontenc}
\usepackage[utf8]{inputenc}
\usepackage{lmodern,amsmath,amssymb,amsthm,mathtools,booktabs,array,longtable}
\usepackage[margin=1in]{geometry}
\usepackage{microtype}
\usepackage{xurl}
\usepackage[hidelinks]{hyperref}
\usepackage{enumitem}
\usepackage{fancyhdr}
\setlength{\headheight}{14pt}
\pagestyle{fancy}\fancyhf{}\fancyhead[L]{\small Fixed-couple m\'enage rows}\fancyhead[R]{\small Report 221}\fancyfoot[C]{\thepage}
\newtheorem{theorem}{Theorem}[section]
\newtheorem{lemma}[theorem]{Lemma}
\newtheorem{corollary}[theorem]{Corollary}
\newtheorem{proposition}[theorem]{Proposition}
\theoremstyle{remark}\newtheorem{remark}[theorem]{Remark}
\hypersetup{pdftitle={Fixed couple menage rows},pdfauthor={Report 221},pdfsubject={Exact conjectures and uniform all-orders asymptotics}}
\newcommand{\fall}[2]{(#1)_{\underline{#2}}}
\newcommand{\ind}{\mathbf 1}
\newcommand{\TV}{d_{\mathrm{TV}}}
\newcommand{\RR}{\mathcal R}
\newcommand{\NN}{\mathbb N}
\DeclareMathOperator{\nint}{nint}
\title{Fixed couple m\'enage rows\\\large Exact differences, concavity, uniform expansions and inverse enclosures}
\author{Report 221}
\date{4 October 2026}
\begin{document}
\maketitle
\begin{abstract}
We prove the adjacent-difference and row-unimodality conjectures displayed in OEIS A332709. A path-rook polynomial identity expresses every adjacent difference in terms of A127548, with a separate zero at the even central step. An explicit injection proves discrete concavity and determines every equality case, including the four-entry maximum plateau in even rows with $n\ge6$ and strict log-concavity in odd rows. An exact decomposition about the row mean yields an expansion to every fixed order, uniformly over all allowed first images. It separates the boundary layers from the interior, gives the fourth-column expansion A258667 through relative order ten, and yields an exact total-variation identity in every row and its sharp asymptotics. Finally, we distinguish exact-range recovery by a smooth asymptotic model from a threshold inverse, proving a two-candidate threshold enclosure with existential constants. The latter is not a certified numerical algorithm. Classical rook theory, the published fixed-couple factorization, ordinary m\'enage asymptotics and general inverse methods are credited. Historical checking is bounded, not an exhaustive novelty claim. A portable exact-arithmetic implementation and separate numerical diagnostics accompany the proofs.
\end{abstract}

\section{The problem and the conclusions}
For $n\ge3$ and $3\le k\le n$, let $T(n,k)$ count the permutations $\pi$ of $\{1,\ldots,n\}$ such that
\[
 \pi(i)\notin\{i,i+1\pmod n\}\quad(1\le i\le n),\qquad \pi(1)=k.
\]
The residue $0$ is represented by $n$. These are the fixed-first-image m\'enage counts in A332709 \cite{oeis332709}. We use the reindexing
\begin{equation}\label{eq:notation}
 A(n,s)=T(n,s+2),\quad 1\le s\le n-2,\qquad
 h=\left\lfloor\frac{n-1}{2}\right\rfloor,\quad q=\min\{s,n-1-s\}.
\end{equation}
Thus $q$ is the distance in column indices from the nearer endpoint, counting the endpoint as distance one. The sum of a row is the ordinary circular m\'enage number $U_n$. The columns $s=1,2,3,4,5$ are respectively A258664, A258665, A258666, A258667 and A258673. In the seating interpretation their clockwise chair separations are $2s+1$.

The displayed A332709 entry, revision \#33 dated 1 February 2021, still labels row unimodality and the adjacent-difference formula as conjectures. The natural domain of the latter includes $4\le k\le n$: the predecessor must itself be a column. We prove
\[
 T(n,k)-T(n,k-1)=\mathrm{A127548}(n-2k+4)
 \quad(4\le k\le n,\ n\ge2k-3).
\]
The even central difference is zero; it is not obtained by substituting zero into A127548. Shevelev and Moses had already observed the row-unimodality pattern in the published 2016 paper \cite{sm}. Their factorization is the starting point here.

The stronger row-shape conclusions are exact for every $n$. The asymptotic results are fixed-order statements as $n\to\infty$, with the relative row expansion uniform in $s$. The inverse results concern a fixed column $s$, not a column varying arbitrarily with $n$. These quantifiers will be kept separate.

\subsection{A fourth-column consequence}
For $a_n=\mathrm{A258667}(n)=A(n,4)$ when $n\ge6$,
\begin{align}\label{eq:fourth}
 \frac{(n-2)a_n}{U_n}={}&1+\frac2{n^4}+\frac{16}{n^5}+\frac{90}{n^6}
 +\frac{456}{n^7}+\frac{2296}{n^8}\notag\\
 &+\frac{12051}{n^9}+\frac{67193}{n^{10}}+O(n^{-11}).
\end{align}
Thus the fourth-column count and the row average first differ at the fourth relative order. The leading equivalence posted in A258667 and proved in the linked Lean source remains correct \cite{oeis258667,lean}. Equation~\eqref{eq:fourth} supplies column-specific higher-order information while preserving that leading equivalence.

\section{Classical rook machinery and conventions}
\subsection{Path polynomials and the fixed-couple factorization}
Let $R_\ell(z)$ be the matching polynomial of a path with $\ell$ edges:
\begin{equation}\label{eq:path}
 R_\ell(z)=\sum_{j=0}^{\lfloor(\ell+1)/2\rfloor}
 \binom{\ell+1-j}{j}z^j,\qquad \ell\ge0.
\end{equation}
A matching is a set of nonconsecutive edge positions; subtracting the mandatory gaps gives the displayed binomial coefficient. Deleting the final edge, according to whether it is used, gives
\[
 R_0=1,\quad R_1=1+z,\qquad R_\ell=R_{\ell-1}+zR_{\ell-2}\quad(\ell\ge2).
\]
For a polynomial $p$ of degree at most $N$, set
\begin{equation}\label{eq:phi}
 \Phi_N(p)=\sum_j(-1)^j[z^j]p(z)(N-j)!.
\end{equation}
The rook inclusion-exclusion formula counts allowed permutations by applying $\Phi_N$ to the rook polynomial of the forbidden board. Indeed, a choice of $j$ nonattacking forbidden squares fixes $j$ row-column assignments and leaves $(N-j)!$ completions; all other choices are inconsistent and contribute zero.

The forbidden row-column incidence graph for a circular m\'enage permutation is a cycle with $2n$ edges. Fixing $\pi(1)=k$ deletes row vertex $1$ and column vertex $k$. These vertices are nonadjacent when $3\le k\le n$. The remaining graph is the disjoint union of paths with $2k-5$ and $2n-2k+1$ edges. Rook polynomials multiply over components, so
\begin{equation}\label{eq:cofactor}
 T(n,k)=\Phi_{n-1}\bigl(R_{2k-5}R_{2n-2k+1}\bigr),\qquad
 A(n,s)=\Phi_{n-1}\bigl(R_{2s-1}R_{2n-2s-3}\bigr).
\end{equation}
This is the known Shevelev--Moses factorization, equations (30)--(31) of \cite{sm}; the path-polynomial method also occurs in the general prefix framework of Kagey \cite{kagey}. Neither the factorization nor rook inclusion-exclusion is claimed as new.

Swapping the two factors proves reflection:
\begin{equation}\label{eq:sym}
 T(n,k)=T(n,n+3-k),\qquad A(n,s)=A(n,n-1-s).
\end{equation}
Every entry is positive: the cyclic shift $\pi(i)=i+k-1\pmod n$ is admissible and has first image $k$. Partitioning all m\'enage permutations by their first image gives
\begin{equation}\label{eq:sum}\sum_{s=1}^{n-2}A(n,s)=U_n.\end{equation}

\subsection{The auxiliary sequence}
For $m\ge1$, let
\begin{equation}\label{eq:L}
 L_m=\Phi_m(R_{2m-2})
 =\sum_{j=0}^{m-1}(-1)^j\binom{2m-1-j}{j}(m-j)!.
\end{equation}
It counts permutations $\sigma$ of $[m]$ satisfying $\sigma(i)\notin\{i,i+1\}$ for $1\le i<m$, with the last row unrestricted. The forbidden graph is a path with $2m-2$ edges and an isolated row vertex. Consequently
\begin{equation}\label{eq:Lbounds}
 0\le L_m\le m!,\qquad (L_1,L_2,L_3,L_4,L_5)=(1,0,1,4,19).
\end{equation}
Writing $\mathcal F(z)=\sum_{r\ge0}r!z^r$ as a \emph{formal} series, the coefficient of $z^m$ in $\mathcal F(z/(1+z)^2)$ is, for $m\ge1$,
\[
 \sum_{r=1}^m(-1)^{m-r}r!\binom{m+r-1}{m-r}=L_m.
\]
The substitution $j=m-r$ proves the equality. Thus $L_m$ is A127548, whose formal constant term is $L_0=1$ \cite{oeis127548}. No analytic convergence is being asserted.

For generating-function algebra only, use the signed Touchard convention
\begin{equation}\label{eq:signed}U_0=1,\qquad U_1=-1,\qquad U_2=0.\end{equation}
For $n\ge3$ these are the usual positive counts. The known ordinary generating function \cite{oeis179,wm} is
\[
 \mathcal U(z)=\sum_{m\ge0}U_mz^m
 =\frac{1-z}{1+z}\mathcal F\left(\frac{z}{(1+z)^2}\right).
\]
Multiplication by $(1+z)/(1-z)$ gives the exact, classical relation
\begin{equation}\label{eq:LU}
 L_m=U_m+2\sum_{r=0}^{m-1}U_r\quad(m\ge1).
\end{equation}
For $m\ge2$ the small signed terms cancel, and this can be written $L_m=U_m+2\sum_{r=3}^{m-1}U_r$ with the sum empty when appropriate. Using $U_1=0$ in \eqref{eq:LU} is incorrect, already at $m=2$.

\section{Exact adjacent differences}
\begin{lemma}\label{lem:poly}
For integers $a\ge0$ and $b\ge a+2$,
\begin{equation}\label{eq:poly}
 R_{a+2}R_b-R_aR_{b+2}=(-z)^{a+2}R_{b-a-2}.
\end{equation}
\end{lemma}
\begin{proof}
Define $F_0=0$, $F_1=1$ and $F_{j+1}=F_j+zF_{j-1}$. Then $R_r=F_{r+2}$. For $q\ge p\ge0$, let
\[
 W_{p,q}=F_{p+1}F_q-F_pF_{q+1}.
\]
The recurrence gives $W_{p,q}=-zW_{p-1,q-1}$ for $p\ge1$, while $W_{0,q-p}=F_{q-p}$. Hence $W_{p,q}=(-z)^pF_{q-p}$. On using the recurrence once, the left side of \eqref{eq:poly} becomes
\[
 R_{a+1}R_b-R_aR_{b+1}=W_{a+2,b+2}=(-z)^{a+2}R_{b-a-2}.
\]
\end{proof}

\begin{theorem}[All adjacent differences]\label{thm:difference}
For $4\le k\le n$, put $m=n-2k+4$. Then
\begin{equation}\label{eq:all-difference}
 T(n,k)-T(n,k-1)=
 \begin{cases} L_m,&m\ge1,\\0,&m=0,\\-L_{-m},&m\le-1.\end{cases}
\end{equation}
In particular, the adjacent-difference conjecture of A332709 holds on its full natural domain $4\le k\le n$, $n\ge2k-3$.
\end{theorem}
\begin{proof}
First suppose $m\ge1$. In Lemma~\ref{lem:poly} set $a=2k-7$ and $b=2n-2k+1$. Its left side is the difference of the rook polynomials for $T(n,k)$ and $T(n,k-1)$. Since $a+2=2k-5$ is odd and $b-a-2=2m-2$, it equals $-z^{2k-5}R_{2m-2}$.

From \eqref{eq:phi}, $\Phi_N(z^dp)=(-1)^d\Phi_{N-d}(p)$ whenever the degrees permit it. Applying $\Phi_{n-1}$ therefore cancels the minus sign and gives $\Phi_m(R_{2m-2})=L_m$, because $(n-1)-(2k-5)=m$.

If $m=0$, the two entries form the central reflected pair in an even row and are equal by \eqref{eq:sym}. If $m<0$, reflection takes the difference to the negative of the left-half difference at $k'=n+4-k$, whose parameter is $-m>0$. This proves all three cases.
\end{proof}
\begin{remark}
The $m=0$ case is essential. A127548 has $L_0=1$, whereas the center difference is zero. Similarly, $k=3$ has no predecessor column in the triangle. Neither case can be added to the positive-parameter formula by a zero-extension convention.
\end{remark}
For later use, the left-half form is
\begin{equation}\label{eq:A-diff}
 A(n,s+1)-A(n,s)=L_{n-2s-2}\qquad(n\ge2s+3).
\end{equation}

\section{Concavity and all equality cases}
\begin{lemma}\label{lem:injection}
For every $m\ge1$, $L_{m+2}\ge L_m$. Equality holds exactly when $m=1$. Also $L_m>0$ for every $m\ge1$ except $m=2$.
\end{lemma}
\begin{proof}
Let $\mathcal S_m$ be the permutations counted by $L_m$. For $\sigma\in\mathcal S_m$, define $\tau\in\mathcal S_{m+2}$ by
\[
 \tau(i)=\sigma(i)\ (i<m),\quad
 \tau(m)=m+2,\quad \tau(m+1)=\sigma(m),\quad \tau(m+2)=m+1.
\]
These are distinct values and hence a permutation. The old restrictions remain valid. Row $m$ receives $m+2$, row $m+1$ receives a value at most $m$, and the last row is unrestricted. The original permutation is recovered from the first $m-1$ values and $\tau(m+1)$, proving injectivity.

For $m\ge2$, the shift by three on $m+2\ge4$ letters is admissible and has value $1$ at row $m$. Every image of the injection has value $m+2$ there, so this shift lies outside the image. Hence the inequality is strict. Direct computation gives $L_1=L_3=1$. Finally, the shift by two is admissible in $\mathcal S_m$ for $m\ge3$, and the remaining cases are in \eqref{eq:Lbounds}.
\end{proof}

\begin{theorem}[Complete row shape]\label{thm:shape}
Every row of $T$ is positive, palindromic, unimodal and discretely concave:
\begin{equation}\label{eq:concavity}
 T(n,k+1)-2T(n,k)+T(n,k-1)\le0\quad(4\le k\le n-1).
\end{equation}
Its maxima are exactly as follows.
\begin{enumerate}[label=(\roman*)]
\item For odd $n\ge3$, the unique maximum is at $k=(n+3)/2$.
\item For even $n\ge6$, precisely $k=n/2,n/2+1,n/2+2,n/2+3$ are maxima.
\item For $n=4$, both existing entries are maxima.
\end{enumerate}
At internal columns, equality in \eqref{eq:concavity} holds exactly at $k=(n+1)/2,(n+5)/2$ for odd $n$ when those columns are internal, and at $k=n/2+1,n/2+2$ for even $n\ge6$. Every odd row is strictly log-concave at every internal column. For even $n\ge6$, equality in log-concavity holds exactly at the two interior columns of the four-entry plateau; for $n=4$ there are no internal columns.
\end{theorem}
\begin{proof}
As $k$ increases by one, the parameter $m=n-2k+4$ decreases by two. Lemma~\ref{lem:injection} and \eqref{eq:all-difference} show that the successive differences are nonincreasing on each side. At the junction they have the forms
\[
 \begin{array}{ll}
 n\text{ odd}:&\ldots,1,1,-1,-1,\ldots,\\
 n\text{ even}:&\ldots,4,0,0,0,-4,\ldots,
 \end{array}
\]
with indices outside the row omitted. Thus they are nonincreasing across the whole row, proving discrete concavity. They are nonnegative on the left and nonpositive on the right, proving unimodality.

For odd $n$, every step toward the center is positive, the final one being $L_1=1$. For even $n$, exactly the final left step $L_2$, the reflected central step, and its right counterpart vanish. When $n\ge6$ this is a four-entry plateau; all steps outside it are strict. The short rows are checked directly.

A second difference vanishes precisely when successive first differences agree. The sole positive-index equality $L_{m+2}=L_m$ is $L_3=L_1$, giving the two stated odd positions. The even equality positions are the two internal points of the zero-difference run. No other equality is possible by strictness in Lemma~\ref{lem:injection}.

For positive consecutive entries $a,b,c$ with $2b\ge a+c$, one has
\[
 b^2\ge\left(\frac{a+c}{2}\right)^2\ge ac.
\]
Equality requires $a=b=c$. It occurs precisely at the two internal points of the even plateau; the odd concavity equalities have unequal neighbors and therefore strict log-concavity.
\end{proof}
The first rows illustrate the distinct parity patterns:
\[
\begin{array}{c|l}
 n& (A(n,1),\ldots,A(n,n-2))\\\hline
 3&(1)\\4&(1,1)\\5&(4,5,4)\\6&(20,20,20,20)\\
 7&(115,116,117,116,115)\\8&(787,791,791,791,791,787).
\end{array}
\]
The theorem concerns a fixed-$n$ row indexed by position. It is distinct in its variable from log-behavior of ordinary or straight m\'enage sequences indexed by $n$; the limited literature check of that subject is discussed in Section~\ref{sec:sources}.

\section{The exact mean decomposition and a uniform tail}
Define the relative row profile
\begin{equation}\label{eq:relative}
 \RR(n,s)=\frac{(n-2)A(n,s)}{U_n}.
\end{equation}
It has average one over the $n-2$ columns. The following identity is the bridge from exact differences to uniform asymptotics.

\begin{theorem}\label{thm:mean}
For every $n\ge3$ and every allowed $s$,
\begin{align}\label{eq:center}
 A(n,s)&=A(n,h)-\sum_{j=q}^{h-1}L_{n-2j-2},\\
 U_n&=(n-2)A(n,h)-2\sum_{j=1}^{h-1}jL_{n-2j-2},\label{eq:total-center}\\
 \RR(n,s)&=1+\sum_{j=1}^{h-1}
 \left(2j-(n-2)\ind_{\{q\le j\}}\right)\frac{L_{n-2j-2}}{U_n}.\label{eq:mean}
\end{align}
Empty sums apply to $n=3,4$.
\end{theorem}
\begin{proof}
Telescope \eqref{eq:A-diff} from the nearer endpoint side to column $h$, then apply reflection. This gives \eqref{eq:center}. For fixed $j\le h-1$, the deficit term $L_{n-2j-2}$ occurs in precisely the $j$ leftmost and $j$ rightmost columns. These sets are disjoint for either parity of $n$. Summing \eqref{eq:center} proves \eqref{eq:total-center}; elimination of the central value proves \eqref{eq:mean}.
\end{proof}

\begin{lemma}\label{lem:factorial}
For every integer $N\ge2$, $\sum_{r=1}^N r!\le2N!$. The same upper bound therefore holds for a sum restricted to either parity.
\end{lemma}
\begin{proof}
The case $N=2$ is immediate. If the inequality holds at $N$, then the next sum is at most $2N!+(N+1)!\le2(N+1)!$, since $N+1\ge2$.
\end{proof}

\begin{theorem}[Uniform fixed-depth truncation]\label{thm:uniform}
For each fixed integer $J\ge0$, as $n\to\infty$,
\begin{equation}\label{eq:uniform}
 \RR(n,s)=1+\sum_{j=1}^{J}
 \left(2j-(n-2)\ind_{\{q\le j\}}\right)\frac{L_{n-2j-2}}{U_n}
 +O_J(n^{-2J-3}),
\end{equation}
uniformly over all $1\le s\le n-2$.
\end{theorem}
\begin{proof}
For large $n$ we have $h>J$. Each omitted weight in \eqref{eq:mean} has absolute value at most $n$, independently of $s$. The omitted indices are positive integers of one parity no larger than $N=n-2J-4$. By \eqref{eq:Lbounds} and Lemma~\ref{lem:factorial}, the omitted numerator is bounded in absolute value by $2nN!$. The classical equivalent $U_n\sim e^{-2}n!$, also proved by the expansion below, yields
\[
 \frac{2n(n-2J-4)!}{U_n}=O_J(n^{-2J-3}).
\]
All constants are independent of the position.
\end{proof}
This proof does not exchange an infinite asymptotic series with a growing sum. It first truncates an exact finite identity with a uniform factorial bound.

\section{An all orders expansion engine}
\subsection{The classical circular expansion}
Write $\fall{v}{k}=v(v-1)\cdots(v-k+1)$ with $\fall{v}{0}=1$, and set
\begin{equation}\label{eq:SM}
 S_M(v)=\sum_{k=0}^M\frac{(-1)^k}{k!\fall{v-1}{k}}.
\end{equation}
For every fixed $M\ge0$, the classical Kaplansky--Riordan expansion is
\begin{equation}\label{eq:Uasymp}
 U_m=e^{-2}m!\{S_M(m)+O_M(m^{-M-1})\}.
\end{equation}
For completeness, a rigorous remainder follows immediately from the stronger Wyman--Moser nearest-integer formula \cite[eq.~(3.23), p.~475]{wm}:
\begin{equation}\label{eq:round}
 U_m=\nint\left(e^{-2}m!\sum_{k=0}^{m-1}
       \frac{(-1)^k}{k!\fall{m-1}{k}}\right),\qquad m\ge2.
\end{equation}
This formula is used only for $m\ge2$; the signed value $U_1=-1$ is separately defined. If $a_{m,k}=1/(k!\fall{m-1}{k})$, then for $0\le k\le m-2$,
\[
 \frac{a_{m,k+1}}{a_{m,k}}=
 \frac1{(k+1)(m-1-k)}\le\frac1{m-1}.
\]
For fixed $M$ and sufficiently large $m$, the absolute tail after $M$ is at most $a_{m,M+1}/(1-1/(m-1))=O_M(m^{-M-1})$. The normalized rounding error is at most $e^2/(2m!)$, smaller than every fixed inverse power. This proves \eqref{eq:Uasymp} and its remainder, including $M=0$. The expansion and nearest-integer refinement are prior results, not contributions of this report.

\subsection{The auxiliary amplitude and coefficient formula}
Set
\begin{equation}\label{eq:EM}
 E_M(v)=S_M(v)+2\sum_{r=1}^{M}\frac{S_M(v-r)}{\fall vr}.
\end{equation}
Then
\begin{equation}\label{eq:Lasymp}
 L_m=e^{-2}m!\{E_M(m)+O_M(m^{-M-1})\}.
\end{equation}
To see this, keep $U_m$ and the $M$ immediately preceding terms in \eqref{eq:LU}. The remaining unsigned count terms have sum at most a constant times $(m-M-1)!$ by $U_r\le r!$ and Lemma~\ref{lem:factorial}; the finitely many exceptional signed terms cancel or are bounded. Dividing by $m!$ makes this $O_M(m^{-M-1})$. Substituting \eqref{eq:Uasymp} in the finitely many retained terms proves \eqref{eq:Lasymp}. There is no derivative assumption on a remainder.

For fixed $d$ and a desired finite order, the coefficients of $L_{n-d}/U_n$ are consequently obtained by expanding
\begin{equation}\label{eq:ratio-engine}
 \frac{E_M(n-d)}{\fall nd\,S_M(n)},
\end{equation}
with $M$ chosen sufficiently large. Each denominator tends to its nonzero leading value after the indicated power of $n$ is factored out, so ordinary finite rational-series algebra is legitimate.

\begin{theorem}[Position-uniform all-orders expansion]\label{thm:allorders}
For every fixed integer $K\ge0$, there are explicitly computable functions $c_0(q),\ldots,c_K(q)$ such that
\begin{equation}\label{eq:allorders}
 \RR(n,s)=\sum_{r=0}^K c_r(q)n^{-r}+O_K(n^{-K-1})
\end{equation}
uniformly over all allowed positions. Each $c_r(q)$ depends on $q$ through finitely many indicators $\ind_{\{q\le j\}}$. It stabilizes at its interior value whenever $2q+1>r$. For fixed $q$, the first difference from the interior expansion is $-n^{-2q-1}$.
\end{theorem}
\begin{proof}
Choose $J$ with $2J+3\ge K+1$ and apply Theorem~\ref{thm:uniform}. Expand each of the finitely many retained ratios by \eqref{eq:ratio-engine}, far enough to allow multiplication by the weight of size $O(n)$. These weights depend on the position only through bounded indicators, so the resulting remainder remains uniform.

The $j$th ratio starts with $n^{-2j-2}$, since both normalized amplitudes have leading value one. Its central-mean contribution $2jL_{n-2j-2}/U_n$ begins at degree $2j+2$, while its boundary subtraction $(n-2)L_{n-2j-2}/U_n$ begins at degree $2j+1$ with coefficient one. For fixed $q$, the first present boundary subtraction is $j=q$, proving the last two assertions.
\end{proof}
Here ``interior'' names the stabilized coefficient regime for a specified finite order; it does not require $q/n$ to stay away from zero. If $q=q(n)\to\infty$, it reaches that regime at every fixed order.

\subsection{Explicit coefficients and the fourth column}
Let $x=1/n$. Table~\ref{tab:coeffs} lists the coefficients at degrees $3$ through $10$ in $\RR$; the constant is one and the degree-one and degree-two coefficients vanish.
\begin{table}[ht]
\centering\small
\begin{tabular}{c|rrrrrrrr}
\toprule
$q$ & $x^3$ & $x^4$ & $x^5$ & $x^6$ & $x^7$ & $x^8$ & $x^9$ & $x^{10}$\\
\midrule
$1$ & $-1$ & $-4$ & $-12$ & $-33$ & $-90$ & $-246$ & $-622$ & $-1067$\\
$2$ & $0$ & $2$ & $15$ & $75$ & $311$ & $1136$ & $3629$ & $9077$\\
$3$ & $0$ & $2$ & $16$ & $90$ & $455$ & $2268$ & $11583$ & $61099$\\
$4$ & $0$ & $2$ & $16$ & $90$ & $456$ & $2296$ & $12051$ & $67193$\\
$\ge5$ & $0$ & $2$ & $16$ & $90$ & $456$ & $2296$ & $12052$ & $67238$\\
\bottomrule
\end{tabular}
\caption{Exact rational-series coefficients, in this case all integers. Each row has a uniform $O(n^{-11})$ remainder in its indicated position regime.}\label{tab:coeffs}
\end{table}
The supplied code computes these by exact fractions, keeping one extra degree before multiplication by $n-2=x^{-1}-2$. This guard is needed to avoid an incorrect last retained coefficient. It also checks the circular amplitude independently through the exact recurrence
\begin{equation}\label{eq:Urec}
 (n-2)U_n=(n^2-2n)U_{n-1}+nU_{n-2}-4(-1)^n\quad(n\ge3),
\end{equation}
and the line amplitude through $L_n-L_{n-1}=U_n+U_{n-1}$, both with exact rational series. Normalized by $e^{-2}n!$, the forcing in \eqref{eq:Urec} is smaller than every inverse power, giving
\[
 u(n)=u(n-1)+\frac{u(n-2)}{(n-1)(n-2)}
\]
as a formal coefficient relation.

For a fixed fourth column, $q=4$ eventually, so Table~\ref{tab:coeffs} gives \eqref{eq:fourth}. It also shows uniformly for $q\ge2$ that
\begin{equation}\label{eq:interior-leading}
 \RR(n,s)=1+2n^{-4}+O(n^{-5}),
\end{equation}
while at either endpoint
\begin{equation}\label{eq:endpoint}
 \RR(n,1)=1-n^{-3}-4n^{-4}-12n^{-5}+O(n^{-6}).
\end{equation}
The posted A258667 leading equivalent $e^{-2}(n-1)!$ agrees with these results because $U_n/(n-2)\sim e^{-2}(n-1)!$. The linked Lean file's concluding assertion is an equivalence statement, not equality of every coefficient. It was read for scope; its compilation was not rerun for this report.

\section{Exact and sharp distance from the uniform distribution}
Choose a circular m\'enage permutation uniformly, and put $S=\pi(1)-2$. Then $\Pr(S=s)=A(n,s)/U_n$. Let $\nu_n$ be the uniform law on $\{1,\ldots,n-2\}$, and define total variation by half the sum of absolute probability differences. We first establish its exact form for every row.

\begin{lemma}[Auxiliary growth and parity tails]\label{lem:TVgrowth}
For $m\ge3$,
\begin{equation}\label{eq:Lrec}L_m=mL_{m-1}-U_{m-2}.
\end{equation}
For every $m\ge2$, $L_{m+1}\ge mL_m$ and $L_{m+2}\ge m(m+1)L_m$. For every $N\ge1$,
\begin{equation}\label{eq:Ltail}
 \sum_{\substack{1\le j\le N\\j\equiv N\ (\mathrm{mod}\ 2)}}L_j\le2L_N.
\end{equation}
\end{lemma}
\begin{proof}
The formal series $\mathcal F(w)=\sum_{r\ge0}r!w^r$ satisfies
$w^2\mathcal F'(w)=(1-w)\mathcal F(w)-1$, by coefficient comparison. Put $g(z)=z/(1+z)^2$ and $\mathcal L(z)=\mathcal F(g(z))$. Since $g'(z)=(1-z)/(1+z)^3$, the chain rule for formal series yields
\[
 z^2\mathcal L'(z)=\frac{1-z}{1+z}(1+z+z^2)\mathcal L(z)-(1-z^2).
\]
Using $\mathcal U(z)=(1-z)\mathcal L(z)/(1+z)$ and rearranging gives
\[
 \mathcal L(z)-1-z^2\mathcal L'(z)-z\mathcal L(z)
 =-z^2\mathcal U(z)-z^2.
\]
Coefficients of degree $m\ge3$ give \eqref{eq:Lrec}. At $m=2$ the extra term matters: $L_2=2L_1-U_0-1=0$. At $m=3$, the signed $U_1=-1$ gives $L_3=1$.

For $m\ge3$, append the fixed point $\sigma(m)=m$ to a circular m\'enage permutation on $[m-1]$. The first $m-2$ restrictions are unchanged. In row $m-1$ the old restriction excludes $m-1$ and $1$, whereas the new one excludes $m-1$ and $m$; the appended value $m$ was absent from the old image. Thus the result is counted by $L_m$, injectively, proving $U_{m-1}\le L_m$. Apply \eqref{eq:Lrec} at $m+1$ to obtain
\[
 L_{m+1}=(m+1)L_m-U_{m-1}\ge mL_m\qquad(m\ge3).
\]
The $m=2$ case follows from $L_3=1$, $L_2=0$, and iteration gives the two-step inequality.

For $N=1,2,3,4$, the parity sums in \eqref{eq:Ltail} are respectively $1,0,2,4$, so the bound holds. For $N\ge5$, induction and growth give
\[
 \sum_{\substack{j\le N\\j\equiv N\ (2)}}L_j
 \le L_N+2L_{N-2}
 \le L_N\left(1+\frac{2}{(N-2)(N-1)}\right)\le2L_N.
\]
There is no division by $L_2$; the zero case was handled separately.
\end{proof}

\begin{theorem}[Exact sign pattern and total variation]\label{thm:TV}
For $n=3,4,6$ the row law is uniform. For $n=5$ and every $n\ge7$, precisely the two endpoints have mass below $1/(n-2)$, and every interior entry has mass strictly above it. For every $n\ge3$,
\begin{align}\label{eq:TV}
 \TV(\mathcal L(S),\nu_n)
 &=2\left(\frac1{n-2}-\frac{A(n,1)}{U_n}\right)\notag\\
 &=\frac2{n^4}+\frac{12}{n^5}+\frac{48}{n^6}+O(n^{-7})\qquad(n\to\infty).
\end{align}
The distance has an expansion to every fixed order from the endpoint coefficients. Through degree ten, its next coefficients are $162$, $504$, $1500$ and $4244$.
\end{theorem}
\begin{proof}
Let $D(n,s)=(n-2)A(n,s)-U_n$. For $n\ge5$, the exact mean decomposition at $q=2$ reads
\begin{equation}\label{eq:D2}
 D(n,2)=2L_{n-4}-\sum_{j=2}^{h-1}(n-2-2j)L_{n-2j-2}.
\end{equation}
The subtracted coefficients are nonnegative and at most $n-6$. The indices run through one parity up to $n-6$. For $n\ge9$, Lemma~\ref{lem:TVgrowth} gives
\begin{align*}
 D(n,2)&\ge2L_{n-4}-2(n-6)L_{n-6}\\
 &\ge2(n-6)^2L_{n-6}>0,
\end{align*}
since $n-6\ge3$. Formula~\eqref{eq:D2} gives $D(5,2)=2$, $D(7,2)=1$ and $D(8,2)=8$. The remaining rows $n=3,4,6$ are the constant rows already displayed.

Reflection and unimodality make every interior entry of a nonconstant row at least its second entry. All its interior deviations are therefore strictly positive. The deviations sum to zero, and the two endpoint deviations are equal, so both are strictly negative. Total variation equals the total negative mass, proving the exact identity. In the three constant rows both sides vanish; the singleton row $n=3$ is handled separately rather than counting two distinct endpoints.

Finally, expand $2(1-\RR(n,1))/(n-2)$ using Table~\ref{tab:coeffs} and $(n-2)^{-1}=n^{-1}(1-2/n)^{-1}$. This proves the coefficients and all-orders extension.
\end{proof}
The exact sign classification requires no asymptotic onset. The numerical onset limitations for the inverse models below are separate and remain in force.

\section{Fixed column inverse models and exact range recovery}
\subsection{Three inverse notions}
Fix an integer $s\ge1$, and define $a_n=A(n,s)$ for integers $n\ge s+2$. It is useful to distinguish:
\begin{enumerate}[label=(\roman*)]
\item the inverse of an explicitly chosen smooth asymptotic model;
\item recovery of $n$ from an exact range value $y=a_n$;
\item the threshold inverse $N_s(y)=\min\{n\ge s+2:a_n\ge y\}$ for a general real target.
\end{enumerate}
They are different objects. In particular, an asymptotically tiny error in a real inverse does not permit a general-input ceiling formula, because an arbitrarily small error can cross an integer. These distinctions and the generic inverse machinery already occur in the repository's canonical and combinatorial transseries texts \cite{canonical,combinatorial}.

Let $P_{s,M}(t)=\sum_{r=0}^M c_r(s)t^r$, the fixed-column relative polynomial from Theorem~\ref{thm:allorders}. Define the specified model
\begin{equation}\label{eq:model}
 F_{s,M}(x)=\frac{e^{-2}\Gamma(x+1)}{x-2}
 S_M(x)P_{s,M}(1/x)
\end{equation}
on a sufficiently large positive half-line. For $M=0$, both last factors are one. The factors approaching one are positive there, and their logarithmic derivatives are $O(x^{-2})$. The digamma expansion gives
\begin{equation}\label{eq:derivative}
 (\log F_{s,M})'(x)=\log x+O_{s,M}(x^{-1}).
\end{equation}
Thus the model is eventually strictly increasing and unbounded.

\begin{theorem}[Recovery at exact range values]\label{thm:range}
For fixed $s,M$, the unique nearby real solution $x_n$ of $F_{s,M}(x_n)=a_n$ exists for all sufficiently large $n$, and
\begin{equation}\label{eq:range}
 x_n=n+O_{s,M}\left(\frac{n^{-M-1}}{\log n}\right).
\end{equation}
In particular, rounding $x_n$ to the nearest integer recovers $n$ eventually.
\end{theorem}
\begin{proof}
The forward expansions give
\[
 a_n=F_{s,M}(n)\{1+O_{s,M}(n^{-M-1})\},\qquad
 \log a_n-\log F_{s,M}(n)=O_{s,M}(n^{-M-1}).
\]
By \eqref{eq:derivative}, the model's logarithm changes by order $\log n$ between $n$ and either endpoint of $[n-1,n+1]$. These endpoint values bracket $\log a_n$ for large $n$. Continuity and strict monotonicity give the unique root there. The mean value theorem, with derivative bounded below by a positive constant times $\log n$, proves \eqref{eq:range}. Its error is eventually less than one half.
\end{proof}
The onset in this theorem is existential. The model is not an exact or canonical interpolation of the integer sequence. The accompanying floating-point roots are diagnostics, not certified round-off or remainder bounds.

\subsection{A Lambert W starting point}
Stirling's formula and the first forward coefficients give, for each fixed $s$,
\begin{equation}\label{eq:log}
 \log a_n=n(\log n-1)-\frac12\log n+c+\frac{13}{12n}+O_s(n^{-2}),
 \qquad c=\frac12\log(2\pi)-2.
\end{equation}
Indeed $\log\Gamma(n+1)$ contributes $1/(12n)$, $-\log(n-2)$ contributes $2/n$, and $\log S_M(n)$ contributes $-1/n$ when expanded to adequate order. The fixed-column relative correction starts only at order $n^{-3}$ or later.

For $y$ large, let
\begin{equation}\label{eq:Lambert}
 L=\log y-c,\qquad X=\frac{L}{W(L/e)},\qquad \ell=\log X,
\end{equation}
where $W$ is the positive real branch, so $X(\log X-1)=L$. At exact range values $y=a_n$,
\begin{equation}\label{eq:inverse-start}
 n=X+\frac12-\frac{23}{24X\ell}
       +O_s\left(\frac1{X^2\ell}\right).
\end{equation}
To prove this directly, substitute $n=X+\delta$ in the smooth main expression of \eqref{eq:log}. The term $\delta\ell-\tfrac12\ell$ forces $\delta_0=1/2$. The residual at order $1/X$ is
\[
 \frac{\delta_0^2}{2}-\frac{\delta_0}{2}+\frac{13}{12}
 =\frac18-\frac14+\frac{13}{12}=\frac{23}{24}.
\]
It is cancelled by $-23/(24X\ell)$. The remaining residual is $O(X^{-2})$, and division by a derivative of order $\ell$ proves \eqref{eq:inverse-start}. Higher orders follow by finite reversion of the explicit model with a correspondingly accurate forward expansion. This is a specialization of standard inversion, not a new Lambert-$W$ technique.

For a precise comparison, choose $M\ge4$, put $G_M(x)=e^{-2}\Gamma(x+1)S_M(x)/(x-2)$, and compare its inverse with that of $F_{4,M}$. Since $P_{4,M}(1/x)=1+2x^{-4}+O(x^{-5})$, the mean value theorem gives, at a common large target,
\begin{equation}\label{eq:model-shift}
 F_{4,M}^{-1}(y)-G_M^{-1}(y)
 =-\frac{2}{X^4\log X}\{1+o(1)\}.
\end{equation}
The positive fourth-column correction lowers the inverse. This compares the two specified smooth models, not two undefined interpolations.

\section{A two candidate threshold enclosure}
The smooth-model range theorem does not by itself solve the threshold problem. The following enclosure is the appropriate consequence when the remainder constants are retained.

\begin{theorem}\label{thm:threshold}
Fix $s\ge1$ and $M\ge0$, and write $F=F_{s,M}$. There are constants $C>0$ and an integer $N\ge s+2$ such that for every integer $n\ge N$,
\begin{equation}\label{eq:envelope-data}
 |\log a_n-\log F(n)|\le Cn^{-M-1},
\end{equation}
and $F_\pm(x)=F(x)\exp(\pm Cx^{-M-1})$ are strictly increasing for $x\ge N$. If
\begin{equation}\label{eq:y-condition}
 y>F_+(N),\qquad y>\max_{s+2\le n<N}a_n,
\end{equation}
where the second restriction is vacuous for an empty prefix, define
\[
 x_+(y)=F_+^{-1}(y),\qquad x_-(y)=F_-^{-1}(y).
\]
Then the exact threshold satisfies
\begin{equation}\label{eq:threshold}
 \lceil x_+(y)\rceil\le N_s(y)\le\lceil x_-(y)\rceil,
\end{equation}
and
\begin{equation}\label{eq:width}
 0\le x_-(y)-x_+(y)
 =O_{s,M}\left(\frac{X^{-M-1}}{\log X}\right),
\end{equation}
where $X$ is from \eqref{eq:Lambert}. Thus, for sufficiently large $y$, the enclosure contains at most two consecutive integers. If it contains two, one exact comparison at the lower candidate determines the threshold.
\end{theorem}
\begin{proof}
Equation~\eqref{eq:envelope-data} follows from the forward relative remainder and positivity. Increase $N$ so that $F$ and both envelopes are increasing: their logarithmic derivatives are $\log x+O(x^{-1})$ by \eqref{eq:derivative}, since the extra derivative is $\mp C(M+1)x^{-M-2}$.

The sequence itself is eventually strictly increasing. To prove this without differentiating an unknown remainder, write $\log a_n=\log F(n)+r_n$ where $r_n=O(n^{-M-1})$. Then $r_{n+1}-r_n=O(n^{-M-1})=O(n^{-1})$, and integration of \eqref{eq:derivative} gives
\[
 \log a_{n+1}-\log a_n=\log n+O(n^{-1})>0
\]
for large $n$. Increase $N$ again if needed. The growth of $F$ and the remainder also make $a_n$ unbounded, so the threshold exists.

The envelope inequality is $F_-(n)\le a_n\le F_+(n)$ for $n\ge N$. Conditions \eqref{eq:y-condition} exclude earlier indices and ensure that both roots lie above $N$. If $a_n\ge y$, then $F_+(n)\ge y$ and $n\ge x_+(y)$, proving the lower bound. Conversely, the integer $n=\lceil x_-(y)\rceil$ satisfies $F_-(n)\ge y$ and hence $a_n\ge y$, proving the upper bound.

Both roots are asymptotic to $X$, by the leading logarithm $x(\log x-1)$. At $x=x_+$,
\[
 \log F_-(x_+)=\log y-2Cx_+^{-M-1}.
\]
At $x_-$ the same logarithm equals $\log y$. The mean value theorem and a logarithmic derivative of order $\log X$ prove \eqref{eq:width}. Eventually this width is less than one, so its ceiling endpoints are equal or consecutive. In the latter case, testing $a_\ell\ge y$ at the lower integer $\ell$ selects $\ell$ if true and $\ell+1$ otherwise, by the enclosure.
\end{proof}
\paragraph{Scope of certification.}
The constants $C,N$ in this theorem exist by the proof, but no numerical certified values are supplied. The theorem is therefore not an executable certified threshold algorithm. No assertion is made that $N_s(y)=\lceil F^{-1}(y)\rceil$ for every sufficiently large input. A practical certified implementation would need explicit global bounds, finite-prefix handling and rigorous root enclosures; floating-point examples do not provide those ingredients.

\section{Reproducibility and implementation safeguards}\label{sec:repro}
The accompanying ZIP contains this \LaTeX{} source, the PDF, a source ledger with version pins, deterministic exact code, optional high-precision diagnostics and a build script. It contains no copied source-paper PDFs, internal research notes or review reports.

The exact core uses only the Python standard library. Its public functions check types and domains, rejecting booleans, floating indices and invalid columns. In particular, the adjacent-difference function distinguishes the center from $L_0$, and the circular count documents its signed small values. It exposes no routine advertised as a certified inverse.

The verification suite checks the path recurrence and polynomial identity, every row through $n=100$, all first and second differences, complete concavity and log-concavity equality classifications, maxima, the mean decomposition, and the $L/U$ relation. A separate subset-DP permanent computes the original allowed-position cofactors through $n=14$. Exhaustive permutations check the line counts and the injection through $m=8$. Exact fraction arithmetic reproduces the table, its order consistency and first boundary separation, the total-variation coefficients and the inverse coefficient residual. The all-row total-variation identity, its precise sign exceptions and the auxiliary growth and tail bounds are checked as well. A recurrence route cross-checks the forward amplitudes.

Deliberately false values and invalid API calls are negative controls: a wrong signed $U_1$, a wrong center step, the interior coefficient substituted for the fourth-column coefficient, a wrong total-variation constant and invalid inputs must all be detected. Large computed integers can be serialized and parsed in bounded decimal chunks without changing the caller's Python digit limit; a 600-row example is tested with the limit set to 640. Checks use explicit exceptions rather than Python \texttt{assert}, so they remain active with \texttt{python -O}. The normal and optimized exact JSON receipts must agree byte for byte.

Numerical diagnostics use \texttt{mpmath} at 100 decimal digits and write a separate JSON file marked diagnostic only. They compare exact counts with truncated ratios and solve the explicit models at exact range targets. They check numerical convergence and implementation consistency; they are not the proof of any asymptotic remainder or certified inverse onset. The build creates a deterministic ZIP by sorting members and fixing archive timestamps and file modes. Rebuilding the \LaTeX{} PDF also fixes its date-dependent metadata through \texttt{SOURCE\_DATE\_EPOCH} and pdf\TeX{} controls. The README specifies dependencies and commands.

\section{Source scope and historical limitations}\label{sec:sources}
The source ledger is bounded to the material actually checked on 4 October 2026. It is not an exhaustive priority search.
\begin{itemize}[leftmargin=*]
\item A332709's internal entry displayed revision \#33, 1 February 2021, with both conjectures. A site-wide footer date is not an entry revision. The full revision history was not audited.
\item The full published Shevelev--Moses paper was read. Its path-rook factorization and fixed-couple formula are credited; its page 10 motivates the row-unimodality question.
\item Wyman--Moser's original paper, especially printed pages 473--475, supplies the nearest-integer remainder argument and explicitly credits the ordinary expansion to Kaplansky--Riordan. These results are classical.
\item Kagey's full 2023 arXiv version was read, including its prefix counting and Fibonacci rook-polynomial framework. No theorem for the present adjacent differences or row concavity was located there. The 2024 publisher version was identified, but its full text returned an access error and was not independently inspected.
\item Gu and Zhao's 2021 publisher abstract and first-page preview were accessible; the full paper was paywalled. Its abstract concerns log-behavior in the size variable for ordinary and straight m\'enage numbers. This is suggestively distinct from fixed-size row concavity, but full non-overlap with the inaccessible paper is not claimed.
\item A258667's displayed 2026 proof comment and complete linked Lean source were read for mathematical scope. They establish leading equivalence. The Lean proof was not rerun, and the present higher-order corrections do not contradict it.
\item The repository's canonical and combinatorial transseries source files were checked at pinned Git blobs. They already supply general inverse, range-recovery and threshold distinctions. Those methods are credited, not rebranded as new here.
\end{itemize}
Targeted identifier and keyword searches found no earlier resolution of the two displayed A332709 conjectures among the inspected sources. The mathematical proofs here stand independently of that bounded search. We claim the stated theorems, not global historical novelty or a complete clearance of all related literature.

\section{Further questions}
Several natural next steps remain separate from the results proved above.
\begin{enumerate}
\item Make the fixed-column inverse constants $C,N$ explicit, control the finite prefix and combine them with rigorous interval roots. This would turn the threshold enclosure into a certified algorithm.
\item Investigate whether similarly explicit total-variation identities hold for other conditioned permutation classes; the exact two-endpoint mechanism here may fail in more complicated forbidden boards.
\item Extend the path-identity approach to several fixed couples or prescribed prefixes. General rook factorizations are available, but row-shape inequalities for the resulting multidimensional arrays need separate proofs.
\item Investigate stronger positivity or higher-difference properties. Ordinary discrete concavity alone does not imply arbitrary higher-order sign patterns.
\item Complete the literature check of the inaccessible final publisher texts before making any historical-priority assertion.
\end{enumerate}

\begin{thebibliography}{99}
\raggedright
\bibitem{oeis332709} The OEIS Foundation, \emph{A332709}, entry by Peter Kagey; inspected 4 October 2026, internal revision \#33, 1 February 2021. \url{https://oeis.org/A332709/internal}.
\bibitem{oeis127548} The OEIS Foundation, \emph{A127548}, formal generating function and identities; inspected 4 October 2026. \url{https://oeis.org/A127548}.
\bibitem{oeis179} The OEIS Foundation, \emph{A000179}, m\'enage numbers and signed Touchard conventions. \url{https://oeis.org/A000179}.
\bibitem{oeis258667} The OEIS Foundation, \emph{A258667}, fourth fixed-couple column, including the 30 June 2026 leading-equivalence proof comment; inspected 4 October 2026. \url{https://oeis.org/A258667}.
\bibitem{sm} V. Shevelev and P. J. C. Moses, \emph{Alice and Bob Go to Dinner: A Variation on M\'enage}, Integers \textbf{16} (2016), A72. \url{https://emis.de/ft/19616}.
\bibitem{kr} I. Kaplansky and J. Riordan, \emph{The probl\`eme des m\'enages}, Scripta Mathematica \textbf{12} (1946), 113--124. Attribution as recorded by Wyman--Moser; not independently audited here.
\bibitem{wm} M. Wyman and L. Moser, \emph{On the probl\`eme des m\'enages}, Canadian Journal of Mathematics \textbf{10} (1958), 468--480. \url{https://doi.org/10.4153/CJM-1958-045-6}.
\bibitem{kagey} P. Kagey, \emph{Ranking and Unranking Restricted Permutations}, arXiv:2210.17021v2, 9 November 2023. \url{https://arxiv.org/abs/2210.17021v2}. Published in Discrete Applied Mathematics \textbf{355} (2024), 247--261, \url{https://doi.org/10.1016/j.dam.2024.05.010}; final publisher full text not inspected.
\bibitem{gu} Y.-X. Gu and F.-Z. Zhao, \emph{The Log-Behavior of M\'enage Numbers}, Results in Mathematics (2021). \url{https://doi.org/10.1007/s00025-021-01468-5}. Abstract and first-page preview only; full text not inspected.
\bibitem{lean} Google DeepMind, \emph{AlphaProof Nexus A258667 Lean source}, concluding \texttt{target\_theorem\_0}. \url{https://github.com/google-deepmind/alphaproof-nexus-results/blob/main/APNOutputs/OEIS/oeis_A258667_conjecture_0.lean}. Source SHA-256 in the supplement.
\bibitem{canonical} V. Reshetnikov, \emph{Transseries: the polynomial-logarithmic calculus, series reversal at infinity, and inversion}, ProveIt repository, canonical source. Git blob \nolinkurl{26378ed517d4ed2e21add80faacfd8851f1fbd10}; inverse distinctions, range recovery and gamma inversion sections. \url{https://github.com/VladimirReshetnikov/ProveIt}. Exact path and SHA-256 in the supplement.
\bibitem{combinatorial} V. Reshetnikov, \emph{Combinatorial Transseries and Their Inverses}, ProveIt repository source. Git blob \nolinkurl{17647298fa31d61c061d556a3b1e451a228ae20e}; inverse notions and interpolation discussion. \url{https://github.com/VladimirReshetnikov/ProveIt}. Exact path and SHA-256 in the supplement.
\end{thebibliography}
\end{document}
